% coset.tex

%%%%%%%%%%%%%%%
\begin{frame}
  \begin{definition}[Coset (陪集))]
    Suppose that $H \le G$.
    For $a \in G$,
    \[
      aH = \set{ah \mid h \in H}, \quad Ha = \set{ha \mid h \in H},
    \]
    is called the \red{left coset (左陪集)}
    and \red{right coset} of $H$ in $G$, respectively.
  \end{definition}
\end{frame}
%%%%%%%%%%%%%%%

%%%%%%%%%%%%%%%
\begin{frame}
  \[
    S_{3} = \set{(1), (1\; 2), (1\; 3), (2\; 3), (1\; 2\; 3), (1\; 3\; 2)}
  \]

  \[
    H = \set{(1), (1\; 2)} \le S_{3}
  \]

  \pause
  \begin{gather*}
    (1) H = H = (1\; 2) H \le S_{3} \\[6pt]
    (1\; 3) H = \set{(1\; 3), (1\; 2\; 3)} = (1\; 2\; 3)H \\[6pt]
    (2\; 3) H = \set{(2\; 3), (1\; 3\; 2)} = (1\; 3\; 2)H
  \end{gather*}
\end{frame}
%%%%%%%%%%%%%%%

%%%%%%%%%%%%%%%
\begin{frame}
  \[
    S_{3} = \set{(1), (1\; 2), (1\; 3), (2\; 3), (1\; 2\; 3), (1\; 3\; 2)}
  \]

  \[
    A_{3} = \set{(1), (1\; 2\; 3), (1\; 3\; 2)} \le S_{3}
  \]

  \pause
  \begin{gather*}
    (1)A_{3} = (1\;2\;3)A_{3} = (1\;3\;2)A_{3} = A_{3} \le S_{3} \\[6pt]
    (1\;2) A_{3} = (1\;3)A_{3} = (2\;3)A_{3} = \set{(1\;2), (1\;3), (2\;3)}
  \end{gather*}
\end{frame}
%%%%%%%%%%%%%%%

%%%%%%%%%%%%%%%
\begin{frame}{}
  \begin{theorem}
    Suppose that $H \le G$, $a, b \in G$.
    \begin{enumerate}[(1)]
      \item
        \[
          a \in aH
        \]
      \item
        \[
          aH = H \iff a \in H
        \]
      \item
        \[
          aH \le G \iff a \in H
        \]
      \item
        \[
          aH = bH \iff a^{-1}b \in H
        \]
      \item
        \[
          \forall a, b \in G\; (aH = bH) \lor (aH \cap bH = \emptyset)
        \]
      \item
        \[
          |aH| = |H| = |bH|
        \]
    \end{enumerate}
  \end{theorem}
\end{frame}
%%%%%%%%%%%%%%%

%%%%%%%%%%%%%%%
\begin{frame}
  \[
    \purple{\boxed{aH = H \iff a \in H}}
  \]

  \vspace{1em}
  \begin{columns}[t]
    \column{0.50\textwidth}
      \pause
      \[
        \brown{\boxed{aH = H \implies a \in H}}
      \]
      \pause
      \[
        a \in aH = H
      \]
    \column{0.50\textwidth}
      \pause
      \[
        \brown{\boxed{a \in H \implies aH = H}}
      \]
      \pause
      \[
        aH \subseteq H H \triangleq \set{hh' \mid h, h' \in H} = H
      \]
      \pause
      \[
        H = (aa^{-1})H = a(a^{-1}H) \subseteq aH
      \]
  \end{columns}
\end{frame}
%%%%%%%%%%%%%%%

%%%%%%%%%%%%%%%
\begin{frame}
  \[
    \purple{\boxed{aH \le G \iff a \in H}}
  \]

  \vspace{1em}
  \begin{columns}[t]
    \column{0.50\textwidth}
      \pause
      \[
        \brown{\boxed{a \in H \implies aH \le G}}
      \]
      \pause
      \[
        a \in H \implies aH = H \le G
      \]
    \column{0.50\textwidth}
      \pause
      \[
        \brown{\boxed{aH \le G \implies a \in H}}
      \]
      \pause
      \[
        a \in aH \implies a^{-1} \in aH \implies e \in aH
      \]
      \pause
      \[
        \exists h \in H\; e = ah \implies a = h^{-1} \in H
      \]
  \end{columns}
\end{frame}
%%%%%%%%%%%%%%%

%%%%%%%%%%%%%%%
\begin{frame}
  \[
    \purple{\boxed{aH = bH \iff a^{-1}b \in H}}
  \]

  \pause
  \[
    \brown{\boxed{a^{-1}b \in H \iff a^{-1}b H = H}}
  \]

  \pause
  \[
    aH = bH \implies a^{-1}aH = a^{-1}bH \implies a^{-1}bH = H
  \]

  \pause
  \[
    a^{-1}b H = H \implies a(a^{-1}bH) = aH \implies bH = aH
  \]
\end{frame}
%%%%%%%%%%%%%%%

%%%%%%%%%%%%%%%
\begin{frame}
  \[
    \purple{\boxed{\forall a, b \in G\; (aH = bH)
      \lor (aH \cap bH = \emptyset)}}
  \]

  \pause
  \[
    \forall a, b \in G\; (aH \cap bH \neq \emptyset \to aH = bH)
  \]

  \pause
  \vspace{0.30cm}
  \begin{center}
    Take any $g \in aH \cap bH$.
  \end{center}

  \pause
  \[
    \blue{\exists h_{1}, h_{2} \in H\;
      (ah_{1} = g = bh_{2})}
  \]

  \pause
  \[
    aH = a(h_{1}H) = (ah_{1})H = (bh_{2})H = b(h_{2}H) = bH
  \]
\end{frame}
%%%%%%%%%%%%%%%

%%%%%%%%%%%%%%%
\begin{frame}
  \[
    \purple{\boxed{|aH| = |H| = |bH|}}
  \]

  \vspace{0.30cm}
  \pause
  \[
    \sigma: aH \to H, \qquad ah \mapsto h
  \]
\end{frame}
%%%%%%%%%%%%%%%

%%%%%%%%%%%%%%%
\begin{frame}
  \begin{center}
    A \cyan{balanced} \red{partition} of $G$ by its subgraph $H$
  \end{center}

  \fig{width = 0.50\textwidth}{figs/partition}

  \[
    \purple{\boxed{\forall a, b \in G\; (aH = bH)
      \lor (aH \cap bH = \emptyset)}}
  \]
  \[
    \purple{\boxed{|aH| = |H| = |bH|}}
  \]
\end{frame}
%%%%%%%%%%%%%%%

%%%%%%%%%%%%%%%
\begin{frame}
  \begin{definition}[Index ($H$ 在 $G$ 中的指标/指数)]
    \[
      G/H = \set{gH \mid g \in G}
    \]

    \[
      [G : H] \triangleq |G/H|
    \]
  \end{definition}

  \pause
  \vspace{0.50cm}
  \begin{theorem}[Lagrange's Theorem]
    Suppose that $H \le G$. Then
    \[
      |G| = [G : H] \cdot |H|
    \]
  \end{theorem}
\end{frame}
%%%%%%%%%%%%%%%

%%%%%%%%%%%%%%%
\begin{frame}
  \begin{theorem}
    \[
      H \le G \implies |H| \red{\mid} |G|
    \]
  \end{theorem}

  \pause
  \vspace{0.50cm}
  \begin{center}
    There are \red{\it no} subgraphs of order 5, 7, or 8 of a group of order 12.
  \end{center}

  \pause
  \begin{theorem}
    Let $G$ be a finite group, $a \in G$. \\[3pt]
    Then the order of $a$ divides the order of $G$.
  \end{theorem}

  \pause
  \vspace{0.50cm}
  \begin{theorem}
    \begin{itemize}
      \item There are only 2 groups of order 4.
      \item There are only 2 groups of order 6.
    \end{itemize}
  \end{theorem}
\end{frame}
%%%%%%%%%%%%%%%